\section{Conclusion}
\label{sec:conclusion}

\vr{} extends \fr{} with a hard VRA rejection rule that enforces majority-minority
district preservation as a chain invariant.
The three-way step accounting identity ($N_{\mathrm{acc}} + N_{\mathrm{vra}} +
N_{\mathrm{mh}} = S$) enables full audit of chain behaviour, and the inherited
\fr{} seed derivation (\texttt{FR\_FORWARD\_}/\texttt{FR\_REVERSE\_}) makes
every VRA rejection independently verifiable.
On NC, GA, and TX, VRA rejection rates are $15$--$25\%$, with EC distributions
$2.7$--$3.0\%$ above the unconstrained baseline---confirming that the VRA
constraint narrows but does not dominate the accessible plan space.

The correct tool for litigation-grade VRA ensemble analysis is
\texttt{--structure ratio-optimal-vra --weights-override vra-aligned
--search vra-recom}: \vraSection{} generates a compliant starting plan,
and \vr{} samples the distribution over compliant alternatives.

\paragraph{Open questions.}
Three design choices are deferred to future work.
(1) \emph{Per-district thresholds}: the current implementation uses a single
global \texttt{vap\_threshold}; allowing per-district thresholds
(\texttt{HashMap<u32, f64>}) would accommodate cases where different districts
have different legal requirements.
(2) \emph{Incremental VAP tracking}: the current VRA check is $O(|D|)$ per
protected district per step; incremental tracking of district VAP fractions
could reduce this to $O(\text{swapped tracts})$ per step, amortising the
check over the recombination boundary.
(3) \emph{Combined minority columns}: \texttt{HVAP+BVAP} is currently computed
at load time; future work may expose this as a first-class \texttt{minority\_col}
option with its own legal characterisation.
